\chapter{元学习}\label{chap:meta}

\begin{learningbox}
本章将介绍元学习（metalearning）。具体而言，学完本章后，你将能够：
\begin{itemize}[leftmargin=*, nosep]
    \item 识别大语言模型训练流程中所采用的不同学习{\em 层次}。

    \item 理解上下文学习、{\em 少样本学习}与元学习之间的联系。

    \item 构建一个由{\em 内}学习循环和{\em 外}学习循环组成的元学习系统。

    \item 认识移动目标问题：为什么在动态元学习系统中实现{\em 稳定学习}具有挑战性。

\end{itemize}
\end{learningbox}

\section{为什么需要元学习？}
前面已经看到，具有自我反思能力的大语言模型方法十分强大，但要让训练过程收敛，可能并不容易。在智能体大语言模型中，学习发生在多个层次，有时这些层次的学习还会同时进行。在机器学习中，我们把能在多个层次上学习的算法称为元学习算法。元学习通常被描述为{\em 学会学习（learning to learn）}，例如，学习哪一种提示更适合某项任务。元学习是一个强有力的概念，但不同学习层次有时会相互干扰，使学习停滞。

为了理解大语言模型的多层次学习算法为何有时不起作用，我们将更深入地考察元学习。本章将讨论，如何提高动态且具有丰富反馈的元学习方法的稳定性。

元学习一词用于指多层次学习算法。元学习算法学习如何学习，以便快速解决新的、从未见过且模型未接受过相应训练的任务中的样例。

\subsection{深度学习很慢}
元学习在机器学习领域已有很长的历史，例如可参见 \citet{schmidhuber1987evolutionary,hospedales2021meta,huisman2021survey}。它受到人类学习的启发，因为人们认为，人类能够学会如何学习。推动元学习发展的另一个原因是速度。

问题在于，常规深度学习很慢。典型深度学习任务使用的神经网络规模庞大，具有数百万乃至数十亿个参数；训练这些网络需要大型数据集，往往包含数百万个样例。从头训练一个网络需要很长时间，从数小时到数月不等。直到采用专用并行硬件（GPU）之后，深度学习才取得成功 \citep{krizhevsky2012imagenet}。

这时，研究者开始憧憬少样本学习（few-shot learning）：只用少量样例进行训练，就能获得高准确率的方法。

\subsection{实现少样本学习}
在机器学习中，对于某些类型的问题，一些算法比另一些算法更合适。此外，大多数训练算法都有一些可以调整以改善性能的要素，例如学习率、动量、参数初始化方案；对于多层次算法，甚至还包括对底层训练算法的选择。这些要素称为{\em 超参数（hyperparameter）}，它们控制着算法如何学习。我们能否通过学习来确定这些要素呢？

不同类型的问题需要不同的超参数设置。例如，对于较平滑的数据集，较高的学习率或局部搜索算法可能效果很好；而对于峰值突出的数据集，较低的学习率或全局搜索算法效果更好。关键在于，机器学习中最佳的超参数配置取决于需要解决的任务，而且尤为重要的是，这种配置可以通过学习得到。针对不同任务，学习算法的这些超参数，就称为{\em 元学习}。（相应地，模型常规参数的学习现在被称为{\em 基础学习（base learning）}。）

元学习的目标，是构建能够更快学会某些任务的算法。我们不应从头开始使用一种尚未适应该任务的通用算法，而应先让学习算法的超参数与待解决的任务类型相匹配。元学习的做法，是先从一组相关任务中学习如何配置学习算法，从而学会更快地学习某一组任务。

学会了这种适配算法的超参数配置后，我们就有望实现少样本训练：只用少量样例进行训练，仍能在新测试任务的样例上获得高准确率。元学习通过从相关知识出发进行训练，而非从头开始，实现少样本学习。

元学习由一个包围着（基础）训练循环的（超参数）训练循环组成。在经典的非神经网络机器学习中，常见的超参数包括学习率或算法选择，后者会在训练任务上测试不同的学习算法 \citep{brazdil2022metalearning}。在深度学习中，经常使用的超参数是神经网络权重的初始化 \citep{finn2017model}。

在智能体大语言模型中，我们也已经遇到不同层次的学习，下一章还将继续讨论多层次学习。因此，研究元学习文献中的理论与认识，对我们很有帮助。

\section{内循环与外循环}

与所有监督学习方法一样，元学习包括两个阶段：(1) 元{\em 训练}阶段，在多种相关任务上训练；(2) 元{\em 测试}阶段，在一个新的、从未见过的任务上测试，考察它对这一新{\em 任务}的泛化能力。目标是学到一组在多种训练任务上都表现良好的训练算法超参数设置，从而改善泛化能力。

训练模型参数的过程通过一个优化循环实现。在循环中，损失函数用于计算训练样例与当前模型估计值之间的差异。训练算法遍历数据集，每次调整模型参数，以使损失总和最小。

学会学习的概念包含两个循环：用于优化模型参数的内训练循环，以及用于优化算法超参数的外训练循环。内循环在一个任务的{\em 样例}上优化参数，外循环则在某个应用领域的{\em 任务集合}上优化超参数。

每个基础学习任务 $\mathcal{T}_i$ 都有一个数据集 $D_i$ 和一个学习目标，即损失函数 $\mathcal{L}_i$。因此，$\mathcal{T}_i=(D_i,
\mathcal{L}_i)$。每个数据集由输入样例与标签的配对组成，即 $D_i=\{(x_j,y_j)\}$，并划分为训练集与测试集：$D_i=\{D_{i,train},D_{i,test}\}$。在每个训练数据集上，依据损失函数 $\mathcal{L}_i(\theta_i,D_{i,{train}})$，估计一个参数化模型 $\hat{f}_{\theta_i}(D_{i,{train}})$。模型 $\hat{f}$ 由一个算法进行近似，该算法受一组超参数 $\omega$ 控制。具体的超参数因算法而异，但常见的超参数包括学习率 $\alpha$、初始参数 $\theta_0$ 以及算法常数。

内循环优化基础参数 $\theta$，外循环优化超参数 $\omega$，以便利用相应的测试数据，在基础任务集合上获得最佳性能：
$$ \omega^\star=\underbrace{\argmin_\omega
  \mathcal{L}}_{\text{外循环}} (\underbrace{\argmin_{\theta_i}
\mathcal{L}_\omega (\theta_i,D_{i,train})}_{\text{内循环}}) .$$
内循环在任务 $\mathcal{T}_i$ 的数据集 $D_i$ 内优化 $\theta_i$，外循环则跨任务和数据集优化 $\omega$。元优化的结果是一组最优超参数 $\omega^\star$。

元学习完成后，我们就得到一组超参数，即一个经过调优的算法，它在某个应用领域的一组任务上表现良好。我们可以假定，只要在一个新任务的少量样例上进行训练，它就能很好地泛化，在未见过的样例上测试时获得高准确率。

\paragraph{迁移学习}
元学习与迁移学习之间存在重要联系。迁移学习先在单个任务上预训练神经网络，再将权重迁移到一个快速微调阶段。因此，迁移学习可以视为一种基于网络初始化的元学习方法的单次迭代：迁移学习就是只有一个基础学习任务的元学习。

\section{大语言模型与元学习}
在大语言模型领域，许多地方都能见到元学习。一种常见的元学习方法，是用先前学到的知识初始化模型参数，使其能够在相关任务上更快地训练。这就是迁移学习，在每一种标准大语言模型训练流程的预训练／微调阶段都能见到：我们取预训练得到的参数，针对特定任务进一步微调。元学习是反复进行的微调，而大语言模型正是元学习训练流程的典型例子，它实现了机器学习的一大目标：少样本学习。上下文学习是机器学习的重要突破之一 \citep{brown2020language}。

进入智能体大语言模型领域后，少样本学习还出现在许多其他地方。我们将讨论静态与动态两种形式的元学习。在静态元学习中，内循环中的权重不受其他循环影响；在动态元学习中，则会受到影响。

\subsection{静态大语言模型元学习}
我们先回顾静态大语言模型元学习的一些例子。

\subsubsection{预训练流程}
常规大语言模型训练流程，是一种复杂的、多阶段迁移学习方案的典型例子。它具有良好的零样本性能，也就是说，它能够根据预训练和微调获得的知识，准确回答许多常规问题，例如：{\em Mozart 出生在哪里？}

训练流程中的不同阶段构成内优化循环，而选择迁移哪些权重，以及人为选择在后续阶段如何微调，则可以视为外循环的超参数选择。

这是一种静态元学习，因为某一阶段的权重一旦训练完成，就不会受到后续阶段训练结果的影响。训练流程的各阶段之间，权重保持稳定。整个流程是一个纯粹顺序进行的训练过程。

\subsubsection{上下文学习}
超越零样本学习，上下文学习是一种少样本学习：用户提供少量训练样例，提示大语言模型通过推理给出更好的答案。预训练模型是内循环，提示中的样例引导外循环。

这是一种静态元学习，因为模型的权重值不受提示中样例的影响。

\subsubsection{思维链}
常规上下文学习将{\em 事实}作为少样本示例，侧重于内容。思维链将{\em 过程步骤}作为少样本示例，侧重于过程。

同样，预训练模型是内循环，提示中的样例引导外循环，而模型的权重保持静态。

\subsubsection{静态元学习是稳定的}
常规大语言模型训练流程由顺序运行的迁移学习阶段，即内循环组成。一旦某个内训练循环收敛，就不再运行该阶段。这种流程是一种静态元学习。（微调结果不会再次反馈给预训练器，预训练器也不会根据微调结果调整超参数后重新运行。不同训练阶段不会相互干扰。）

与静态元学习相对，还存在动态元学习，其中不同迭代确实会相互干扰。

\subsection{动态大语言模型元学习}
智能体大语言模型使用了不同形式的动态元学习。在动态元学习中，迭代并发运行，可能使彼此的收敛过程变得不稳定。

前面几章讨论了自我反思与 RLVR。在这些方法中，外循环可以改变内循环，一个内循环的结果也可能影响另一个内循环。这样，我们就创造了一个移动的训练目标。这可能使训练过程不稳定并发散。训练结果也可能表现得十分脆弱：输入值或超参数的微小变化，例如不同的随机种子、提示的不同改写方式，甚至元学习过程本身时长的变化，都可能引起很大的反应。此外，当同一模型被用于不同角色，例如纠错与生成时，基础大语言模型产生幻觉的倾向可能被放大。

下面讨论智能体大语言模型中的动态元学习例子。

\subsubsection{自我反思}
在自我反思中（第~\ref{sec:selfreflection}~节），智能体控制框架负责提示管理，以不同角色调用大语言模型，例如评估一个状态或改进一个答案。

用元学习的术语来说，{\em 模型}是内训练循环，而{\em 控制框架}是外循环。虽然模型本身是静态的，但它会由不同提示调用，控制框架又会利用模型的回答产生新提示。因此，提示所产生的结果（外循环）开始相互干扰，而且可能以非线性、非直观的方式交互。

此外，当同一模型被用于不同角色时，其缺陷，例如偏差和幻觉，可能被放大。

\subsubsection{RLVR}
可验证奖励强化学习（reinforcement learning with verifiable rewards，RLVR，第~\ref{sec:rlvr}~节）是编程与数学应用中一种流行的训练方案，它将推理阶段的推理结果用于训练阶段的微调。

用元学习的术语来说，{\em 微调}是内循环，而推理阶段的{\em 推理}是外循环。RLVR 在内循环中使用外循环的结果，将推理结果反馈给训练阶段，混合了两个学习过程。应当注意，避免那些被当作真值使用的推理阶段错误，在微调／推理循环中被放大。

\section{稳定动态元学习}
当内、外循环之间的干扰处于欠阻尼状态时，就会出现不稳定的训练。此时，内、外循环之间的振荡可能不断增大，而非逐渐消失。

可以通过增大阻尼，或将不同内循环发生变化的时间尺度解耦，来缓解这一问题。例如，可以降低某个循环的学习率，或者让某些内循环的参数保持稳定的时间长于其他循环。另一个思路是混合使用不同类型或规模的模型，例如混合专家（Mixture of Experts）\citep{jiang2024mixtralexperts}。学生／教师设置会有意让较强与较弱的模型分别承担不同角色。通过使用不同模型，我们有望抑制幻觉，以及内循环之间的相互干扰。

一般而言，当元学习表现异常时，我们应当尝试增加多样性，促进探索，防止贪心行为导致局部最优或训练发散。这些解决办法往往需要大量实验，才能逐渐建立对问题的直觉。调试智能体大语言模型可能很有挑战性。

\section{结语}
正如前面所见，我们可以将大语言模型训练理解为一种元学习。常规大语言模型训练是一种迁移学习，而智能体大语言模型所采用的推理阶段推理与训练，则运用了更多元学习思想。虽然自我反思等动态元学习方法的理念十分强大，但仍需谨慎，因为内、外循环之间的相互干扰可能使训练过程变得不稳定。这种情况可能意外地出现在此前一直运行正常的智能体中。

当我们在一个与世界交互的智能体中使用训练好的模型，而且还利用这种交互的反馈训练该智能体时，就可以把模型视为基础学习器，即内循环，把与世界的交互视为元学习器，即外循环。前面关于训练流程内元学习的所有注意事项，也适用于训练流程外的元学习。只是此时与现实世界发生交互，错误的后果可能更加严重。

在线机器学习的相关领域，例如深度强化学习，同样面临训练不稳定和移动目标问题。深度强化学习领域已经提出许多稳定化技术，例如经验回放缓冲区和目标网络 \citep{mnih2015human,hessel2018rainbow}。这些方法通常试图将两个学习过程的周期解耦，往往通过减慢其中一个过程来实现。层次推理模型（Hierarchical Reasoning Models，HRM）\citep{wang2025hierarchical} 也引入了类似方法。HRM 采用双层元学习方法，同样通过以两种不同速度学习来稳定学习过程。

设计智能体大语言模型解决方案时，了解这些稳定化方法很有帮助。下一部将概述智能体大语言模型与世界交互的不同方式。这些方式从静态逐步过渡到动态。

\bibliographystyle{plainnat}
\bibliography{agenticbook,refs,case,synthetic}
